\hypertarget{comparing-four-ways-to-assist-diagram-construction}{%
\section{Comparing four ways to assist diagram
construction}\label{comparing-four-ways-to-assist-diagram-construction}}

The useful comparison is between the operations an assistant performs
and the evidence each operation adds. An interactive editor helps an
analyst see and revise a hypothesis. A language model can propose
mechanisms and missing variables from descriptions. A response-based
conditional-independence procedure can organize model judgments into a
graph search. A workflow agent can select and operate a statistical
discovery method under declared assumptions. These operations may share
an interface, but they answer different scientific questions. The
proposed product should preserve that distinction when it combines them.

CausalChat illustrates the interactive route. Its inspected preprint
describes an optional CPDAG obtained from data, GPT-4 comments on
relations, visual summaries linked to model responses, and
structural-equation coefficient updates after graph changes. This is a
useful design for making an analyst's reasoning inspectable. Its
evidence for interface benefits nevertheless needs careful
interpretation. Six university students first resolved three
self-selected edges using BIC feedback, then another three using LLM
text, and finally the remaining three using CausalChat Lite. The
reported error counts were 5/18, 7/18 and 1/18. The edge sets and order
changed together with the interface. Those results describe the authors'
experiment; they do not identify the effect of replacing one interface
with another on matched, randomly assigned banking tasks. Expert
participants also questioned missing citations and the quality of the
scientific evidence behind model suggestions. A compact visualization
may improve comprehension while leaving the causal premise unresolved.
\citep{v12_la_2410_14146}

The breadth-first method of Jiralerspong and colleagues offers a
different contribution. A query expands one variable into its proposed
effects, and insertion rejects an edge that would create a directed
cycle. The authors' linear query-count argument follows from visiting
each variable once. It does not make the total input/output token count
linear: each expansion can refer to many variables and the accumulated
graph. The reported 221-node experiment achieved an F1 score of 0.351,
while some comparison configurations and the proposed variant using
observational statistics were omitted for resource reasons. That
evidence motivates an efficiency baseline with actual tokens, latency
and omitted requests accounted for. It supplies no reason to certify
every retained arrow merely because the final graph is acyclic.
\citep{v12_la_2402_01207}

chatPC places language-model judgments inside a conditional-independence
search. Its prompts include variable descriptions, context and a
conditioning set; its repeated-answer test concerns the relative
probabilities of YES and NO responses. The paper reports agreement under
exchange of the two queried variables of 0.510 for GPT-3.5 and 0.870 for
GPT-4, with the latter often rejecting independence. Symmetry is a
necessary property of a conditional-independence relation. Measuring the
frequency with which a model respects it is useful, but that frequency
is not the error rate of a statistical test in the bank's observation
law. The two sources of randomness must remain separate: a model can
consistently repeat a mistaken judgment about the world.
\citep{v12_la_2406_07378}

The causal-learn+ system description instead places agents around formal
discovery: data inspection, assumption elicitation, method
recommendation, tool coordination and interpretation. It separates
suggestions, user decisions, algorithm settings and graph outputs, and
emphasizes retaining CPDAG and PAG ambiguity. This is a strong
architectural alternative when the aim is access to an existing formal
method. Its call for future benchmarks of diagnosis, selection and
calibrated interpretation also limits the performance conclusions
available from the description itself. Keeping a formal algorithm at the
center does not establish its premises, and user approval does not turn
a disputed premise into empirical truth. Conversely, a fully specified
and evaluated hybrid method is conceivable; the architecture's
restrictive preference is not an impossibility theorem about every use
of language-model information. \citep{v12_la_2606_23608}

These comparisons support a product with several explicitly named
operations. Proposing a missing common cause should return a hypothesis
and its evidence. Editing a reviewed graph should return a new graph and
the changed premises. Running a specified discovery method should retain
the method's actual output class. Identifying a target should consume
that class through a supported formal path. A user can benefit from all
four services without being led to believe that a plausible graph
proposal has already passed the last two.

\hypertarget{a-measured-risk-score-is-not-the-risk-variable-it-names}{%
\section{A measured risk score is not the risk variable it
names}\label{a-measured-risk-score-is-not-the-risk-variable-it-names}}

Variable elicitation is consequential before the first edge is drawn.
Consider a stipulated servicing example. Let \(U\) be a customer's
unobserved financial resilience, \(A\) an assistance offer, \(Y\)
successful repayment, and \(Z\) a pre-offer score category. Take all
variables to be binary. The generated mechanisms are

\[
 P(U=1)=\tfrac12,\qquad
 P(A=1\mid U=u)=\tfrac15+\tfrac35u,
\]

\[
 P(Y=1\mid A=a,U=u)=\tfrac1{10}+\tfrac15a+\tfrac12u,
 \qquad P(Z=u\mid U=u)=\tfrac45.
\]

The score error is independent of the action and outcome mechanisms
conditional on \(U\). The offer is more frequent in the more resilient
category in this stipulated assignment law. The graph contains
\(U\to A\), \(U\to Y\), \(U\to Z\) and \(A\to Y\). A dictionary that
calls both \(U\) and \(Z\) ``customer risk'' can conceal the measurement
distinction without changing a single recorded number.

Under an intervention on \(A\), the distribution of \(U\) remains one
half in each category. The mean success probabilities are

\[
 E\{Y(0)\}=\tfrac7{20},\qquad
 E\{Y(1)\}=\tfrac{11}{20},\qquad
 \tau=\tfrac15.
\]

In the observed groups, however, \(P(U=1\mid A=1)=4/5\) and
\(P(U=1\mid A=0)=1/5\). Their observed difference is therefore

\[
 E(Y\mid A=1)-E(Y\mid A=0)
 =\tfrac15+\tfrac12\left(\tfrac45-\tfrac15\right)=\tfrac12.
\]

Adjustment for the noisy score improves neither the scientific premise
by declaration nor the graph by relabeling. For \(Z=1\), Bayes' rule
gives

\[
 P(U=1\mid A=1,Z=1)=\tfrac{16}{17},\qquad
 P(U=1\mid A=0,Z=1)=\tfrac12.
\]

For \(Z=0\), the corresponding probabilities are \(1/2\) and \(1/17\).
The residual difference is \(15/34\) in either score stratum, so the
standardized contrast is

\[
 \sum_z\{E(Y\mid A=1,z)-E(Y\mid A=0,z)\}P(z)
 =\tfrac15+\tfrac12\tfrac{15}{34}
 =\tfrac{143}{340}\simeq .421.
\]

It remains above the true effect by \(15/68\), despite using a variable
measured before the offer. All sixteen cells of the joint generated law
of \((U,A,Z,Y)\) have positive probability. This is not a failure caused
merely by an empty observed stratum. It is the consequence of treating a
proxy as if it blocked the latent common-cause path. Exact enumeration
reproduces each calculation. \citep{exec13_diagram_examples}

An assistant can help by asking what the score measures, who constructed
it, what data were available at scoring time and whether the assignment
policy used the score or additional information. It can also propose
alternative diagrams: the offer may depend only on the recorded score,
or it may depend on latent information also seen by a servicing officer.
Those alternatives have different adjustment implications. A
language-model statement that the score is ``a confounder'' is too
coarse to settle the question.

For a proposal, publicly available dictionaries and controlled examples
are sufficient to demonstrate this failure and specify the required
interface. They cannot establish the target bank's true measurement
process. The immediate design should retain the conceptual variable,
measured proxy and measurement mechanism separately. An eventual bank
study of the proxy is justified only if the prototype and decision
analysis show that the unresolved measurement information is
consequential.

\hypertarget{what-additional-measurement-information-would-buy}{%
\section{What additional measurement information would
buy}\label{what-additional-measurement-information-would-buy}}

The proxy example also gives a concrete route beyond a generic request
for better data. Suppose a symmetric error probability \(e\) is known,
and the stronger nondifferential measurement condition

\[
 P(Z\mid U,A,Y)=P(Z\mid U)
\]

holds. For fixed \((a,y)\), write \(o_z=P(A=a,Y=y,Z=z)\) and
\(q_u=P(A=a,Y=y,U=u)\). Then

\[
 \begin{pmatrix}o_0\\o_1\end{pmatrix}
 =
 \begin{pmatrix}1-e&e\\e&1-e\end{pmatrix}
 \begin{pmatrix}q_0\\q_1\end{pmatrix}.
\]

The determinant is \(1-2e\). For the declared range \(0\leq e<1/2\),
inversion yields

\[
 q_u=\frac{(1-e)o_u-eo_{1-u}}{1-2e}.
\]

Thus known measurement information and the stated mechanism permit
reconstruction of the latent joint law in this finite example.
Adjustment using that reconstructed \(U\) law recovers \(1/5\), exactly
matching the intervention calculation. No observed customer identity or
latent label has been manufactured: the reconstruction depends
explicitly on an assumed, known error matrix.
\citep{exec13_diagram_examples}

The route has a meaningful failure boundary. In the vector one-norm, the
inverse matrix has norm \(1/(1-2e)\), because each column's absolute
entries sum to that value. At \(e=1/5\), a worst-case perturbation in
the observed joint probabilities can be amplified by \(5/3\) before the
further ratios needed for adjustment. As \(e\) approaches one half,
measurement inversion becomes unstable. At one half, the score supplies
no information about which latent category generated it, and this
inversion is unavailable. An exact symbolic inverse alone does not make
the resulting effect precise.

In finite samples the inverted entries can become negative. Such entries
may reflect sampling error, a misspecified measurement model or an error
matrix incompatible with the observed law. Clipping them to zero and
renormalizing changes the estimator and its interpretation. A defensible
alternative is to retain a constrained uncertainty region for the
observed probabilities and measurement parameters, then project the
compatible latent laws into effect values. Whether that region has
justified coverage is an additional statistical obligation; the exact
population example does not supply it.

If \(e\) is estimated from a labeled subset, the subset's selection and
label truth matter. A convenience sample of easily resolved customers
need not identify the same error matrix for difficult cases. If
difficulty affects the error rate even after conditioning on \(U\), the
symmetric matrix is wrong. If the score was constructed using post-offer
outcomes, nondifferential measurement can fail. The assistant's useful
contribution is to identify these propositions, locate relevant
documentation and propose a sensitivity analysis. Repeating a confidence
rating about the score cannot replace the missing information.

This makes the next experiment more specific. A public or generated
benchmark can compare inversion, naive proxy adjustment and a valid
refusal as the measurement error changes. It can establish that the
proposed software preserves the declared assumptions and numerical
limits. A later target-bank measurement audit would answer a different
question: whether an adequate and sufficiently representative
measurement model exists there. The expense of that audit belongs after
evidence that the affected decision warrants it.

\hypertarget{post-action-engagement-can-be-a-collider}{%
\section{Post-action engagement can be a
collider}\label{post-action-engagement-can-be-a-collider}}

The label ``engagement'' often invites adjustment, because it sounds
like a useful customer characteristic. Its clock and causal role matter
more than its name. Consider another exact example with independent fair
binary variables \(A\) and \(U\). Let \(A\) be a randomized outreach
offer, \(U\) latent willingness to repay, \(Y=U\), and \(M\) a
post-offer engagement indicator satisfying

\[
 P(M=1\mid A=a,U=u)=\tfrac1{10}+\tfrac25a+\tfrac25u.
\]

The graph has \(A\to M\leftarrow U\to Y\). The offer has zero causal
effect on repayment, and the unadjusted observed difference is zero. All
eight observed cells of \((A,M,Y)\) have positive probability.
Conditioning on \(M\) nevertheless makes the randomized action
informative about \(U\).

Among engaged customers,

\[
 P(Y=1\mid A=1,M=1)=\frac{.9}{.5+.9}=\tfrac9{14},
 \qquad
 P(Y=1\mid A=0,M=1)=\frac{.5}{.1+.5}=\tfrac56.
\]

Their difference is \(-4/21\). Among customers with \(M=0\), the
corresponding probabilities are \(1/6\) and \(5/14\), giving the same
difference. Standardizing over \(M\) therefore reports
\(-4/21\simeq-.190\) even though the true effect is zero. The apparently
richer adjustment set has created the discrepancy.
\citep{exec13_diagram_examples}

The example separates three questions that an assistant should not
merge. The total effect concerns \(E\{Y(1)-Y(0)\}\). An observed
comparison among engaged customers conditions on \(M=1\) after the
action and can compare different latent populations across arms. A
principal-stratum question instead refers to a joint potential-response
condition, such as customers who would engage under either offer
assignment. Its definition and identifying assumptions are different
again. A clearer English description does not make those targets
interchangeable.

If engagement is a mediator in another plausible graph, holding it fixed
can remove part of the pathway through which the offer works. A
controlled direct effect would involve \(Y(a,m)\) and an intervention on
engagement, if that intervention has a coherent meaning. Ordinary
regression adjustment for an observed mediator does not automatically
identify that effect, especially with hidden causes of both the mediator
and repayment. The assistant should first expose the target choice and
then ask which scientific conditions support the requested calculation.

The remedy is not a universal ban on using any variable recorded after
the offer. A genuinely baseline quantity may be entered into a database
later, and retrospective measurement can be relevant to an analysis if
its measurement and selection assumptions are justified. Conversely, a
variable entered early may encode a forecast produced using information
unavailable to the policy at the decision time. The product needs
distinct event, measurement and availability clocks, together with an
account of how the variable was constructed.

\hypertarget{time-unrolling-and-the-difference-between-analysis-and-deployment}{%
\section{Time unrolling and the difference between analysis and
deployment}\label{time-unrolling-and-the-difference-between-analysis-and-deployment}}

An institution may describe a feedback loop: offers affect balances, and
balances affect later offers. A graph over two timeless variables,
``offer'' and ``balance,'' can then contain a directed cycle. That
observation does not establish that every causal model of the workflow
requires a cyclic graph. With dates made explicit, a possible structure
is

\[
 A_t\longrightarrow B_{t+1}\longrightarrow A_{t+1}
 \longrightarrow B_{t+2}.
\]

The arrows describe distinct variable instances. An acyclic
time-unrolled model can represent feedback across decisions while
retaining a clear within-date ordering. Whether it is scientifically
adequate still depends on the sampling frequency, contemporaneous
mechanisms, hidden information and the policy that chooses each action.
Deleting an inconvenient arrow from the timeless graph would conceal
these issues rather than resolve them.

For each instance, the proposed variable record distinguishes the period
described, the event that generated its value, the time it was first
available and any subsequent correction. A customer's end-of-month
balance may be known only after settlement. A revised macroeconomic
release can describe an earlier month while containing information that
was not available when an offer was chosen. Using the revised value to
reconstruct a historical policy can create information leakage without
introducing an obviously backward arrow into a diagram.

Write \(\mathcal H_t\) for the information actually available when an
action is selected. A deployable policy has the form
\(A_t=\pi_t(\mathcal H_t)\), perhaps with explicit randomized variation.
A retrospective adjustment variable need not belong to exactly the same
information set: it might help describe pre-action confounding after
being measured later. The scientific validity of that retrospective use
and the feasibility of using it in a prospective policy are separate
questions. The diagram assistant should preserve both records rather
than infer feasibility from a successful estimation call.

A bank also needs to distinguish a true event from the observation of
that event. An offer can affect whether a customer responds to a survey,
whether a missing identifier is resolved or whether an account enters
the analysis extract. If \(S\) denotes observation, an effect among
\(S=1\) can concern a selected population whose composition changes with
the action. Adding an ``observed'' flag to a data table is not enough;
its parents belong in the causal model. Public synthetic benchmarks can
vary this mechanism deliberately and reveal whether the interface
silently drops the selection node.

These distinctions guide the assistant's questions. A temporal
contradiction should generate a request to refine the variable instance
or clock. A change in data availability should affect the deployable
policy record. An unresolved observation mechanism should remain a
possible source of selection bias. None requires the assistant to invent
a hidden construct or declare a fully specified longitudinal
identification result before the required evidence exists.

\hypertarget{evidence-for-an-arrow-a-path-and-an-exclusion}{%
\section{Evidence for an arrow, a path and an
exclusion}\label{evidence-for-an-arrow-a-path-and-an-exclusion}}

Expert statements differ in logical strength. ``The offer precedes the
withdrawal'' can rule out a particular backward-in-time mechanism under
the declared clock. It does not establish an arrow from offer to
withdrawal, exclude an earlier common cause or prove a nonzero effect.
``The offer changes the fee, which changes withdrawal'' proposes a
directed path through a mediator; it does not necessarily propose a
direct arrow after that mediator is included. ``The offer was
randomized'' concerns an assignment mechanism and its scope. The
assistant should preserve these differences before converting prose into
software constraints.

The same caution applies to exclusions. Absence of an association in one
study can reflect imprecision, cancellation, measurement error or a
different population. Absence of a retrieved passage can reflect search
failure. Neither observation automatically justifies a forbidden causal
path. A hard exclusion removes models from the admitted class; it is
therefore a scientific restriction with consequences for every
downstream identified set.

Let \(\mathcal M\) be the currently admitted set of causal models and
\(K\) a proposed restriction. The restricted class is

\[
 \mathcal M_K=\{m\in\mathcal M:m\text{ satisfies }K\}.
\]

For a fixed observed law and target, its compatible effect set is
contained in the corresponding unrestricted set. Narrowing follows from
set inclusion, not from new empirical accuracy. If the true model is
excluded by \(K\), the narrower result can be more decisive and less
defensible at the same time. If the restriction leaves no compatible
model, the output should identify the contradiction rather than
substitute an empty effect with the meaning ``no effect.''

A useful evidence record therefore includes the exact proposition,
variable instances, population and period, primary source passage,
whether the assertion is observed or assumed, and counterevidence. It
also records the target calculations that depend on the proposition. If
two experts disagree about a common cause, the analyst should be able to
see the effect bounds and feasible actions under each branch. If they
disagree about a remote edge that changes neither the relevant
functional nor the feasible decision, the next review can focus
elsewhere without declaring the remote edge resolved.

One source can support several claims without providing several
independent pieces of evidence. A guideline, a review quoting that
guideline and three generated summaries of the review may all trace to
one original study. The proposed evidence graph should retain that
common origin. It should also preserve differences between a source's
measured population and the bank's target population. Authenticity,
applicability and inferential strength are separate properties of the
same document.

\hypertarget{agreement-between-agents-does-not-measure-causal-certainty}{%
\section{Agreement between agents does not measure causal
certainty}\label{agreement-between-agents-does-not-measure-causal-certainty}}

Consensus can be useful for finding inconsistent explanations, but its
interpretation depends on how the answers were generated. Let
\(V_1,\ldots,V_n\) be binary endorsements of one proposed edge. In a
stipulated exchangeable response model, each has mean \(q\), variance
\(q(1-q)\), and common pairwise correlation \(\rho\). Expanding the
variance of the mean gives

\[
 \operatorname{Var}(\bar V)
 =\frac{q(1-q)}{n}\{1+(n-1)\rho\}.
\]

The variance-equivalent count of independent responses is

\[
 n_{\rm eff}=\frac{n}{1+(n-1)\rho}.
\]

For ten responses and an assumed correlation of .9, this count is
\(100/91\), barely more than one. This is an exact implication of the
response model, not an estimated correlation or a calibration result for
any deployed agent. Shared training, retrieved passages or an initial
proposed answer can create dependence whose size must be studied rather
than guessed. \citep{exec13_diagram_examples}

Even \(\rho=0\) would not solve the more fundamental problem. The
parameter \(q\) describes the chance that the prompted system endorses
an edge. It is not the probability that the edge is true in the target
bank. Independent repetitions can estimate a consistently wrong
endorsement probability very precisely. Assigning different role names
to the same model changes the prompts; it does not establish either
statistical independence or independent access to scientific evidence.

The hallucination-reduction study by Sng and colleagues makes a related
distinction concrete. Its contradiction metric compares model ratings,
its retrieval experiment supplies a mocked true causal graph, and its
debate setup uses opposing positions and an arbiter. The authors
explicitly distinguish reduced contradiction from accurate causal
prediction and leave true-graph validation to further work. That is
useful evidence about a proposed consistency remedy, with a narrower
claim than the paper title might suggest. \citep{exec13_hallucination}

The proposed product can use debate to expose an omitted assumption, ask
for a counterexample or compare the scope of two sources. It should
retain the actual reasons and original passages. A final arbiter's vote
should remain a workflow decision or an explicitly uncalibrated
preference. If the intended downstream operation needs probabilities
over models, those probabilities require a declared prior/likelihood
construction or a separately justified calibration target. A voting
fraction cannot silently supply them.

\hypertarget{reuse-decisions-for-three-available-package-paths}{%
\section{Reuse decisions for three available package
paths}\label{reuse-decisions-for-three-available-package-paths}}

The inspected pywhy-llm revision provides useful hypothesis-elicitation
components. Its identification suggester delegates a back-door
suggestion to the model's confounder suggester, while the inspected
front-door method is a TODO with a pass body and the longer
estimand-construction block is commented out. Its MIT license and
declared Python/LLM dependencies make reuse of selected interaction
ideas plausible. They do not make a returned confounder list an
identified adjustment set. The proposed integration should accept
suggestions as candidate evidence records and send the resulting
reviewed graph to a formal identification operation. It should not
present the front-door placeholder as an available scientific
capability. \citep{v12_la_package_pywhy_llm}

The pinned causal\_GPT source is valuable as an experimental comparison
for response-based discovery. It aggregates model answers through
voting, weighted voting and a proportion test, then maps the selected
decision to zero or one in the conditional-independence interface. Those
endpoints are software signals; they are not calibrated p-values for the
bank's joint distribution. The source also reuses YES and NO counts from
the same answer collection, which makes their covariance relevant even
to the response-level test. The analysis in
Chapter\textasciitilde{}\ref{ch:la03-discovery} develops that separate
issue. The appropriate decision is to retain the method as a labeled
benchmark candidate, with its BSD three-clause obligations and actual
dependencies, while declining direct use as a calibrated sampling-law
oracle. \citep{v12_la_package_causal_gpt}

CausalAgent offers an end-to-end application and MCP comparison. At the
inspected revision, its server advertises PC and OLC tools. It imports
the PC analysis function, while the OLC function called in that module
is neither imported nor defined there; the surrounding exception handler
returns a JSON failure string. A separate output helper takes the text
between the first and last braces and returns an empty dictionary when
decoding fails. These static findings constrain adoption of the
inspected path. They are not a claim that every revision or every
deployment behaves identically. The proposal should reuse the idea of
explicit tool handoffs, require typed invalid/unavailable outcomes, and
defer the OLC path until its concrete implementation is repaired and
actually exercised. \citep{v12_la_package_causalagent}

There are three distinct reuse decisions here. Selected elicitation
components can be integrated behind a proposal-only boundary. An
imperfect scientific method can remain a useful experimental comparator.
A concrete incomplete tool path can be rejected for direct adoption
while retaining other architectural lessons. Porting an entire
application or dependency stack does not resolve any of the scientific
defects. A port should be considered only after a selected path has a
justified purpose and a measured constraint that the port could address.

Actual deployment would still require frozen package and model
identities, license notices, data-handling decisions, runtime budgets,
endpoint conversion checks and independent reference cases. The current
comparison is based on saved source and primary text. No live-model
accuracy, end-to-end latency or runtime conformance is inferred from
that inspection. These limitations define the first prototype trials
without blocking the present architecture and literature argument.

\hypertarget{editing-a-graph-without-losing-its-scientific-history}{%
\section{Editing a graph without losing its scientific
history}\label{editing-a-graph-without-losing-its-scientific-history}}

A graph edit should change one declared object and make its consequences
inspectable. Adding a latent node is different from adding an observed
column. Resolving a circle endpoint is different from moving a node on
the screen. Accepting an expert exclusion is different from an algorithm
finding an orientation under a registered test. The interface should
retain the origin and logical status of each change even when it
displays the resulting graph in one view.

Suppose the analyst begins with a partially directed graph and a
disputed assertion that a risk score is sufficient for adjustment. A
useful edit sequence first records the score as a measured proxy, then
creates a branch with the latent common cause, and finally compares the
target functional under the original and revised assumptions. The proxy
example shows why the changed value can be scientifically meaningful
even when model fit remains good. The analyst should be able to return
to the earlier branch and see the exact premise responsible for the
difference.

Constraint checking also needs a clear purpose. Rejecting cycles
maintains the declared DAG representation, but greedy rejection can
depend on insertion order. Given proposed arrows \(A\to B\), \(B\to C\)
and \(C\to A\), accepting the first two rejects the third; another order
rejects a different arrow. The algorithm has selected one internally
consistent subset without establishing which arrow was scientifically
wrong. The remedy is to expose the conflicting set, preserve the
evidence for each member and seek a revised clock, representation or
premise. A deterministic insertion order improves reproducibility; it
does not supply a causal reason for the choice.

The downstream dependency record should bind an identified expression to
the graph, query and assumptions that produced it. If an edit changes
the outcome's definition, the target changes. If an edit changes a
common-cause restriction, identification may change. If only a displayed
label changes through a valid bijection, the scientific object may
remain the same. These cases deserve different invalidation rules.
Reusing an old estimate because the edited diagram still has the same
number of edges is not a defensible rule.

Review should preserve disagreement when available evidence cannot
resolve it. A graph collection is not automatically a statistical
confidence set, and a union of its effect candidates is not
automatically exhaustive over every causal model. The product must state
which alternatives were considered, why others were excluded and what
additional uncertainty remains. A well explained partial answer can
still narrow the next experiment or rule out a premature investment.

\hypertarget{which-disagreement-is-worth-resolving}{%
\section{Which disagreement is worth
resolving}\label{which-disagreement-is-worth-resolving}}

The monetary question can organize expert review without pretending to
settle causal truth. Consider an assumed proposal scenario in which two
admitted model branches imply effect values .02 and .08 on repayment
probability. Assign an assumed USD 1,000 value to an additional
successful repayment and a USD 50 offer cost. The corresponding net
values per account are

\[
 v_L=1000(.02)-50=-30,\qquad
 v_H=1000(.08)-50=30.
\]

With no justified probability over the two branches, the action is not
robustly beneficial. Relative to choosing the better action in each
branch, both adopting and declining have maximum regret USD 30 per
account. The calculation locates a consequential disagreement; it does
not choose a probability because a majority of prompted agents favored
one branch. \citep{exec13_diagram_examples}

If a sponsor explicitly adopts a planning probability \(\pi\) for the
high-value branch, the expected net value of adoption is \(60\pi-30\).
The best current expected value is \(\max(0,60\pi-30)\). Perfect
information about the branch would allow adoption only when its value is
positive, yielding expected value \(30\pi\). The expected value of
perfect information is therefore

\[
 \operatorname{EVPI}(\pi)
 =30\pi-\max(0,60\pi-30).
\]

At \(\pi=1/2\), it is USD 15 per account; at \(\pi=4/5\), it is USD 6.
For 10,000 otherwise comparable accounts with additive values and one
shared model uncertainty, those assumed upper bounds are USD 150,000 and
USD 60,000. The higher prior confidence reduces the incremental value of
learning even though it makes immediate adoption more attractive. These
are exact planning calculations, not observed bank returns or elicited
institutional beliefs. \citep{exec13_diagram_examples}

A real experiment provides imperfect information, so its expected value
of sample information cannot exceed the perfect-information value under
the same correct decision model. Delay, implementation, review and
measurement costs reduce its net value further. An experiment costing
more than this upper bound is not justified by this particular decision
alone. A lower cost is necessary but not sufficient: one must still show
that the experiment can distinguish the relevant branches and that its
result can change an available action.

A public benchmark can establish that the prototype preserves branch
identities, computes the right functional in a known mechanism and
avoids the proxy/collider mistakes. It generally cannot reveal which
causal mechanism operates in a specific bank. The proposal should
therefore sequence the investment: first test the accessible
computational and reasoning obligations; then ask whether the remaining
target-specific uncertainty is valuable enough to investigate. Expensive
authentic evidence is a conditional next investment, not a prerequisite
for writing a defensible proposal.

\hypertarget{a-benchmark-for-diagram-assistance-with-explicit-failure-credit}{%
\section{A benchmark for diagram assistance, with explicit failure
credit}\label{a-benchmark-for-diagram-assistance-with-explicit-failure-credit}}

Evaluation should begin with a declared task object. One case might ask
for a diagram consistent with a timing description; another for possible
omitted variables; a third for an intervention answer over unresolved
graph alternatives. They need different references. An author-stipulated
graph can provide controlled truth for a generated case. A public
real-world dataset may provide an observation law and metadata without a
unique causal graph. Scoring both against one arbitrarily selected edge
list would confuse missing information with model error.

The proxy and collider cases provide exact formal baselines. The
reference is the stated structural law and its separately enumerated
intervention distribution. The assistant sees the permitted variable
descriptions and observational information, with the hidden generator
truth withheld whenever the task concerns discovery. A correct output
should distinguish what follows from those observations from what
depends on an additional assumption. A refusal to declare
unconfoundedness is useful when the evidence does not supply it; a
refusal to do any arithmetic in a fully specified exact model is a
different failure.

The comparison should include a deterministic dictionary-and-clock
validator, an unassisted formal method in its admitted domain, a
single-model elicitation route, a retrieval-assisted route and a
reviewed human workflow where feasible. Public or generated cases permit
these designs without requiring a bank deployment. A live-model
comparison additionally needs actual model execution, recorded settings
and measured resource use. An offline schema or literature synthesis
supplies neither live accuracy nor superiority.

Perturbations should test meanings that ought to remain invariant.
Permuting variable order should not change the meaning of conditional
independence. Renaming variables through a bijection should preserve a
formal answer when their descriptions remain equivalent. Conversely,
changing a variable from pre-offer state to post-offer response should
change the admissible reasoning. Removing a cited premise should weaken
the claims that depended on it. A benchmark that rewards invariance to
every wording change would miss this distinction between irrelevant
phrasing and substantively changed information.

Scores should retain unsupported, malformed and unavailable requests in
the denominator. They should distinguish valid abstention from an
uninformative missing answer, graph syntax from target identification,
and a justified ambiguous effect set from an omitted calculation. The
monetary example supplies an additional loss measure: did the assistant
recommend an action that the admitted models do not justify, or identify
a potentially valuable next information request? Edge accuracy remains
useful, but it is insufficient when a small number of consequential
edges determines the decision.

Finally, reference information and retrieval material must be separated
deliberately. Giving a model the true graph can test whether it uses
that information correctly, as in a controlled consistency task. It
cannot also measure discovery of that graph from evidence that did not
contain it. The proposed evaluation should report which service was
actually tested, preserve unfavorable cases, and measure the full cost
of retrieving, checking and acting on an answer. That produces a
reviewable reason to build or narrow a diagram-assistance prototype
while leaving target-bank scientific truth to the later study that could
establish it. \label{exec13:diagram-development-end}
